\section{Four Color via the Algebraic $\mathbb{F}_2$ / Jordan Route}
\label{sec:goertzel-4cp-state}

This section is a research-program lens on Goertzel's algebraic route to the
Four Color Theorem (the v23/v24 lineage). It is written to point creative
effort, not to claim a proof. Every conjecture below is labelled. The Lean
identifiers cited are live surfaces in the \texttt{GraphTheory/FourColor}
development; the verbal claims about what is proved versus open track that tree.

\subsection{The creative bet}

The route reframes four-colorability as a single statement of \emph{algebraic
control over a chain space} $\mathbb{F}_2$. Color edges with a Tait edge
coloring (the nonzero elements of $\mathbb{F}_2^2$), read every chain as an
$\mathbb{F}_2$-vector on edges, and let path-xor (symmetric difference) be the
vector addition. In this language the obstruction to a proper coloring lives in
the \emph{boundary-zero Kirchhoff subspace}
\[
  W_0(H) \;=\; \verb|theorem49BoundaryZeroKirchhoffSubspace|,
\]
the intersection of the cycles vanishing on the selected boundary
(\verb|planarBoundaryZeroSubmodule|) with the Kirchhoff (vertex-degree-zero)
relations at the purified interior vertices (\verb|kirchhoffSubmodule|).

The bet is that this subspace is exactly \emph{spanned and detected} by the
projections of Kempe-closure generators (Definition 4.8,
\verb|kempeClosureGeneratorFamily|, projected via
\verb|projectedKempeClosureGeneratorSubspace|). If a single
\emph{boundary-zero projected-generator detector} works --- every nonzero
boundary-zero chain pairs nontrivially, under the $\mathbb{F}_2$ form
\verb|chainDot|, with some projected generator --- then the synthesis endpoint
\verb|theorem49BoundaryRootSynthesis| (manuscript Theorem~4.9) closes and the
coloring follows. The genius of the framing is that it converts a global
topological fact (planarity) into a finite linear-algebra detection question:
the cycle space is small, the generator family is explicit, and the pairing is
computable.

What is genuinely built and machine-checked:
\begin{itemize}\itemsep1pt
  \item the F$_2$ path-xor / boundary-zero detector machinery
    (\verb|chainDot|, \verb|BoundaryZeroProjectedGeneratorDetector|),
  \item the boundary-zero Kirchhoff subspace and its purified interior-vertex
    discharge (\verb|theorem49BoundaryZeroKirchhoffSubspace|,
    \verb|selectedBoundaryInteriorEdgeEndpointVertices|),
  \item the exact v23 Step~2 single-component purification
    \(X_f^{ab}(C)\oplus X_f^{ab}(C^D)=c\cdot 1_{D\cap\partial f}\)
    (\verb|v23_single_component_purification|) and its sum to the
    boundary-only generator,
  \item the edge-cardinality / finrank bound that turns a rank deficit into a
    forced nonzero chain
    (\verb|theorem49BoundaryZeroKirchhoffScalarConstraintSpace_finrank_eq|,
    \verb|exists_nonzero_f2ScalarChain_mem_ker_of_finrank_lt_card|),
  \item the cross-component bridge counterexample and the
    \emph{detector-strength} theorem that $W_0(H)$ is a \emph{proper} subspace
    of the full boundary-zero space on every shell-bearing embedding
    (\verb|theorem49BoundaryZeroKirchhoffSubspace_lt_planarBoundaryZeroSubmodule_of_hasUnblockedInteriorEndpoint|).
\end{itemize}

The last point is the sobering one and the generative one at once. ``Evaders''
--- nonzero boundary-zero chains invisible to every \emph{forced} edge --- exist
because the Kirchhoff target sits strictly inside the cycle space. So the
detector cannot be a cheap dimension count; it must be a genuine claim about the
geometry of Kempe components. That is precisely where the open crux lives.

\subsection{Possible avenues still open}

\paragraph{A1. The explicit-family detector ladder.} The v24 work proved that
the synthesis depends only on the underlying edge-Kempe closure class, and that
\emph{any} detecting coloring family inside the base closure already upgrades to
the full Theorem~4.9 endpoint
(\verb|theorem49BoundaryRootSynthesis_of_familyPairingSeparates| and the
coordinate/edge-predicate variants). This means the live obligation no longer
has to be stated on the whole closure. The avenue: scale finite detector
certificates by \emph{neighborhood enlargement} (radius-$0$, radius-$1$, $\dots$)
rather than by enumerating the closure. The two minimal certificates
(wheel radius-$0$, shared-interior-pair radius-$1$) are already kernel-checked;
the next step is a parametrized radius-$k$ certificate generator.

\paragraph{A2. A computable $\mathbb{F}_2$ elimination layer.} The honest v24
status note flags this as the unblocking move: build a Lean-computable finite
search over the chain space so the benchmark discharge widens safely beyond the
wheel/shared family. Avenue: a verified Gaussian-elimination kernel over
\verb|F2| edge-chains, certified against \verb|chainDot| nondegeneracy, that
either exhibits an evader or certifies its absence on a given embedding.

\paragraph{A3. Residual / live-boundary peeling (v23.5).} The ambient IDBRAD
terminal-peeling interface is \emph{provably impossible} on honest closed-walk
geometry with interior support
(\verb|not_nonempty_interiorDualBoundaryRootAdjDistancePeelData_of_interiorEdgeSupport_nonempty_of_faceBoundaryClosedWalkGeometry|).
The repair is to chase the \emph{moving} frontier $B_k$ of the surviving
residual $\mathcal{R}_k$ rather than the frozen ambient annulus boundary. Avenue:
formalize the residual equation object and prove a live-boundary terminal
theorem. Caution: Lean already shows the present residual wrapper is
\emph{equivalent} to the collar-layer and height-ordered surfaces, so a naive
re-encoding will be rejected by the regression guard. The avenue is only live if
the live-boundary certificate is genuinely new.

\subsection{Conjectured workarounds for the known blocks}

The two named blocks are (i)~the algebraic-cancellation oracle and (ii)~the
planar/Jordan cyclic-separator (Sublemma~6.8). Both are stated below as targets
for genius, not as results.

\begin{conjecture}[Speculative --- Local Kempe detector]
\label{conj:4cp-local-detector}
There is a constant $k$ (independent of the graph) such that on every
shell-bearing planar embedding, the radius-$k$ projected-generator detector
already separates every nonzero boundary-zero Kirchhoff chain: no evader
survives within distance $k$ of a forced edge. If true, the
\verb|BoundaryZeroProjectedGeneratorDetector| reduces to a \emph{finite local}
check, and A1's radius-$k$ certificate generator closes the oracle uniformly.
\end{conjecture}

\noindent\emph{Why it might hold.} The detector-strength theorem says the gap
$W_0(H)\subsetneq$ cycle space is nonempty exactly when an interior endpoint is
unblocked; that endpoint is a \emph{local} feature. The conjecture asserts the
\emph{cause} of every evader is itself local, hence detectable at bounded
radius. \emph{Why it might fail.} Evaders could be globally entangled across the
embedding; a single bad cross-component bridge already breaks the naive count.
The conjecture is the bet that planarity forbids such global entanglement of
boundary-zero chains.

\begin{conjecture}[Speculative --- Jordan separator as an $\mathbb{F}_2$ duality]
\label{conj:4cp-jordan}
Sublemma~6.8 (a planar/Jordan cyclic-separator extracting the annulus
decomposition from boundary-order data) can be replaced by a purely algebraic
$\mathbb{F}_2$ statement: the boundary-zero subspace admits an orthogonal
splitting, under \verb|chainDotBilinForm|, into an ``inside'' and ``outside''
block whenever the embedding carries an honest closed-walk source. The Jordan
curve theorem then enters only through the \emph{nondegeneracy} of the form, not
through an explicit topological separator.
\end{conjecture}

\noindent\emph{Why it might hold.} The honest closed-walk source already encodes
the cyclic order needed for separation; \verb|chainDotBilinForm| is proved
nondegenerate, so an orthogonal Hodge-style splitting of the cycle space is
plausible. \emph{Why it might fail.} The current development shows the unordered
face-incidence API is strictly weaker than ordered/cyclic/closed-walk data; the
algebraic splitting may secretly re-import the very ordered data it claims to
avoid. The conjecture is responsible only if the splitting is derived from
\verb|FaceBoundaryClosedWalkGeometry| and never from a smuggled ambient order.

\begin{conjecture}[Speculative --- Purification entry point]
\label{conj:4cp-purification}
The exact Step~2 purification
(\verb|v23_single_component_purification|, summed to the boundary-only
generator) is the \emph{correct} algebraic entry point: given
Conjecture~\ref{conj:4cp-local-detector}, the boundary-only generator already
spans $W_0(H)$ on every shell-bearing embedding, so no ambient peel theorem is
needed at all. The blocked IDBRAD lane was solving a problem the local detector
makes vanish.
\end{conjecture}

\subsection{Where to point creative effort next}

The single highest-leverage target is
\textbf{Conjecture~\ref{conj:4cp-local-detector} via avenue A1+A2}: a
parametrized radius-$k$ detector certificate generator backed by a verified
computable $\mathbb{F}_2$ elimination layer. Everything else in the route is
either proved (the algebra, the purification, the cardinality bound) or is a
re-statement that the regression guard will reject. The whole edifice rests on
one question --- \emph{do evaders exist at unbounded radius, or can a local
Kempe detector see them all?} --- and that question is now phrased finitely
enough to attack by machine search before betting a human proof on it. The
recommended creative push: enlarge the certificate radius on adversarially
chosen shell-bearing benchmarks (wheel-with-inner-triangle, shared-interior-pair,
cross-component bridge) and watch whether the detector ever needs unbounded
neighborhoods. If it does not, Conjecture~\ref{conj:4cp-local-detector} becomes
a credible human-proof target; if it does, the evader at the boundary of
detection is itself the new mathematical object worth naming.

\subsection*{PLN lens}

\begin{table}[h]
\centering
\small
\begin{tabular}{lcccc}
\toprule
Central node & STV $\langle s,c\rangle$ & ITV $[L,U]$ & Progress & Status \\
\midrule
Algebraic-cancellation oracle (detector) & $\langle 0.30, 0.55\rangle$ & $[0.17,\,0.62]$ & --- & open crux \\
Jordan / cyclic-separator (Sublemma 6.8) & $\langle 0.30, 0.45\rangle$ & $[0.14,\,0.69]$ & --- & open crux \\
F$_2$ algebra + purification + finrank   & $\langle 0.95, 0.85\rangle$ & $[0.81,\,0.96]$ & --- & built \\
\midrule
\textbf{Route (through)} & $\langle 0.25, 0.66\rangle$ & $[0.165,\,0.505]$ & $\sim$55\% & \\
\bottomrule
\end{tabular}
\caption{PLN summary for the 4CP algebraic route. $s$ = expectation the route
closes; $c$ = evidential confidence; indefinite interval $[L,U]=[s\cdot c,\,
s\cdot c+(1-c)]$. Progress is route-exploration completeness. The built-algebra
node is high; the two open cruxes carry the residual risk, so the through-route
expectation $0.25$ sits below the algebra node and reflects the unresolved
detector/separator questions.}
\end{table}

\noindent Through-route: $\langle 0.25, 0.66\rangle$, ITV
$[0.25\cdot 0.66,\;0.25\cdot 0.66+(1-0.66)] = [0.165,\,0.505]$,
PROGRESS $\sim 55\%$. The relatively high progress reflects how much of the
algebraic spine is machine-checked; the modest strength reflects that both
remaining cruxes are genuinely open and the detector-strength theorem proves the
target is strictly nontrivial on every shell.
